\documentclass[11pt]{article}
\usepackage[utf8]{inputenc}
\usepackage[margin=1in]{geometry}
\usepackage{enumitem}
\usepackage{hyperref}
\usepackage{bookmark}
\usepackage{booktabs}
\usepackage{graphicx} 
\usepackage{xcolor}
\usepackage{listings}
\usepackage{float}
\usepackage{tabularray}

\graphicspath{{image/}{./}}  

\definecolor{codegreen}{rgb}{0,0.6,0}
\definecolor{codegray}{rgb}{0.5,0.5,0.5}
\definecolor{codepurple}{rgb}{0.58,0,0.82}
\definecolor{backcolour}{RGB}{220,255,255}  
\definecolor{keywordblue}{RGB}{0,0,180}
\definecolor{commentgray}{RGB}{100,100,100}

\lstset{
    basicstyle=\ttfamily\small,
    numberstyle=\tiny\color{codegray},
    numbers=left,
    numbersep=5pt,
    breakatwhitespace=false,
    breaklines=true,
    captionpos=b,
    keepspaces=true,
    showspaces=false,
    showstringspaces=false,
    showtabs=false,
    tabsize=2,
    frame=single,
    rulecolor=\color{gray!30},
    frameround=tttt,
    backgroundcolor=\color{backcolour},
    xleftmargin=10pt,
    xrightmargin=10pt,
    aboveskip=10pt,
    belowskip=10pt,
}

\lstdefinelanguage{Python}{
    keywords={False, None, True, and, as, assert, async, await,
              break, class, continue, def, del, elif, else, except,
              finally, for, from, global, if, import, in, is, lambda,
              nonlocal, not, or, pass, raise, return, try, while,
              with, yield},
    morekeywords={print, len, range, type, int, str, float, list,
                  dict, tuple, set, open, close, read, write},
    keywordstyle=\color{keywordblue},
    comment=[l]{\#},
    morecomment=[s]{'''}{'''},
    morecomment=[s]{"""}{"""},
    commentstyle=\color{commentgray}\itshape,
    stringstyle=\color{codepurple},
    morestring=[b]{"},
    morestring=[b]{'},
    sensitive=true,
    alsoletter={\_,\@},
}

\lstdefinelanguage{RISCV}{
  morekeywords=[1]{add, sub, and, or, slt, addi, andi, ori, slti, bne, jal, lw, sw},
  morekeywords=[2]{x0, x1, x2, x3, x4, x5, x6, x7, x8, x9,
                   x10, x11, x12, x13, x14, x15, x16, x17, x18, x19,
                   x20, x21, x22, x23, x24, x25, x26, x27, x28, x29,
                   x30, x31,
                    },
  % 注释
  morecomment=[l]{\#},
  morecomment=[s]{/*}{*/},
  % 字符串
  morestring=[b]",
  morestring=[b]',
  sensitive=true,
  alsoletter={.,\_},
}

\lstdefinestyle{mypython}{
    language=Python,
    keywordstyle=\color{keywordblue},
    commentstyle=\color{codegreen}\itshape,
    stringstyle=\color{codepurple},
    backgroundcolor=\color{backcolour},
}


\lstdefinestyle{myriscv}{
    language=RISCV,
    keywordstyle=[1]\color{keywordblue},
    keywordstyle=[2]\color{black},
    commentstyle=\color{codegreen}\itshape,
    stringstyle=\color{codepurple},
    backgroundcolor=\color{backcolour},
}

\lstnewenvironment{python}[1][]{
    \lstset{style=mypython,#1}
}{}

\lstnewenvironment{riscv}[1][]{
    \lstset{style=myriscv,#1}
}{}

\title{Report on miniCPU Project}
\author{Mingxuan Li [2025533130] \quad Ruiyang Xu [2025531125] \quad Kaiwen Yang [2025531126]}
\date{}

\begin{document}

\maketitle

\section{Task Description}
miniCPU Design and Implementation is a cutting-edge educational initiative focused on computer architecture.

In this project, our work revolves around building the datapath based on \texttt{Logisim}, developing an assembler in \texttt{Python}, and writing test cases.

\section{Team's Timeline}

\subsection{Week 1}
In week 1, we gained preliminary understanding of the concept of electrical level and sequential logic, and have studied the structure of registers. We also preliminarily developed a Python-based assembler.

\subsection{Week 2}
In week 2, we studied the fundamentals of computer architecture, including learning RISC-V instruction set and gaining an understanding of the datapath of a single-cycle CPU. Simultaneously, we refined our assembler, enabling it to read assembly files in the terminal, convert assembly instructions into hexadecimal machine code, and output the results via a text file.

\subsection{Week 3}
In week 3, we implemented a single-cycle CPU in Logisim capable of executing 13 RV32I instructions: \texttt{add, addi, sub, and, andi, or, ori, slt, slti, lw, sw, bne} and \texttt{jal}.

\subsection{Week 4}
In week 4, we designed test cases and conducted debugging and modifications on the CPU.

\section{CPU Implementation}
\subsection{Overall Design Plan}
\begin{itemize}
    \item \textbf{Instructions Set}
    \begin{table}[H]
    \centering
    \renewcommand{\arraystretch}{1.5} % Adds vertical padding for a cleaner look
    \begin{tabular}{|c|c|c|c|c|}
        \hline
        \textbf{R-Type} & \textbf{I1-Type} & \textbf{S-Type} & \textbf{J-Type} & \textbf{B-Type} \\ \hline
        add, sub, or, and, slt & addi, andi, ori, lw, slti & sw & jal & bne \\ \hline
    \end{tabular}
    \end{table}
    \item \textbf{Modules Design}
    \begin{itemize}
        \item \textit{Instruction Fetch}: IFU includes the PC, NPC, and IM, where the PC register has an asynchronous reset function with a starting address of 0x00000000.
        \item \textit{Instruction Memory}: IM is a read-only memory module that stores the machine code instructions in hexadecimal format generated by the assembler.
        \item \textit{Register File}: RF contains 32 registers, each 32 bits wide. It supports two asynchronous read ports and one synchronous write port, and has an asynchronous reset function that initializes all registers to 0.
        \item \textit{ALU}: ALU performs arithmetic and logical operations, supporting operations such as addition, subtraction, bitwise \texttt{AND}, bitwise \texttt{OR}, and set less than.
        \item \textit{Data Memory}: DM is a synchronous read-write memory module. It has an asynchronous reset function that initializes all memory locations to 0.
        \item \textit{Extend}: The extend module is responsible for extension of immediate values.
        \item \textit{Control Unit}: The control unit generates control signals for each module and data path based on the opcode and funct3/funct7 fields of the instruction.
    \end{itemize}
    \item \textbf{Overall Structure}
    \begin{figure}[H]
        \centering
        \includegraphics[width=1.0\textwidth]{Overall.png}
        \caption{The overall structure}
        \label{fig:Overall Structure}
    \end{figure}

    \item \textbf{Key Modules}
    \begin{enumerate}
        \item \textit{PC}
        \begin{table}[htbp]
        \centering
        \begin{tblr}{
            colspec = {l X[j]}, % 'l' for left column auto-width, 'X[j]' for flexible justified column
            hlines, % Horizontal lines
            vlines, % Vertical lines
            row{1} = {font=\bfseries}, % Bold header row
            column{1} = {font=\bfseries}, % Monospace font for port names
        }
        \textbf{Port} & \textbf{Description} \\
        clk & Receives the clock signal (rising edge triggered). \\
        reset & Receives an asynchronous reset signal, which sets the current instruction to the starting position when it is $1'b1$. \\
        D & Receives the next instruction address. \\
        Q & Outputs the current instruction address. \\
        \end{tblr}
        \caption{PC Module Ports}
        \end{table}

        \begin{figure}[H]
            \centering
            \includegraphics[width=0.2\textwidth]{PC.png}
            \caption{PC Module Implementation}
            \label{fig:PC Module}
        \end{figure}
        
        \item \textit{Instruction Memory}

        We use a ROM as the instruction memory.  
        
        The physical address space of the instruction memory (i.e., the actual address range connected to the ROM) is 0x0000 to 0xFFFF (a total of 64KB of space). Since the Program Counter (PC) is 32 bits wide, but the instruction ROM only supports a 16-bit address input (a 32-bit address would require excessively large memory and is unnecessary), we only use the lower 16 bits of the PC (PC[15:0]) as the access address for the ROM. Each instruction is 32 bits wide (4 bytes), but the memory is word-addressable (32 bytes per word), where each address location stores 32 bits. Therefore, the address of the next instruction in this ROM is the current address plus 1.

        \begin{figure}[H]
            \centering   % Left Image
            \begin{minipage}{0.25\textwidth}
            \centering
            \includegraphics[width=\textwidth]{Instr_Memory.png}
            \end{minipage}
            \hspace{0.1\textwidth} % Adds horizontal space between the images
            \begin{minipage}{0.25\textwidth}% Right Image
            \centering
            \includegraphics[width=\textwidth]{PC_to_ROM.png}
            \end{minipage}
            \caption{Instruction Memory Addressing Method}
        \end{figure}
    
        \item \textit{Register File}
        \begin{figure}[H]
            \centering
            \includegraphics[width=0.15\textwidth]{Register_file_overall.png}
            \caption{Register File Overall Structure}
            \label{fig:RF Module Overall Structure}
        \end{figure} 

        \begin{table}[htbp]
        \centering
        \begin{tblr}{
            colspec = {l X[j]}, % 'l' for left column auto-width, 'X[j]' for flexible justified column
            hlines, % Horizontal lines
            vlines, % Vertical lines
            row{1} = {font=\bfseries}, % Bold header row
            column{1} = {font=\bfseries}, % Monospace font for port names
        }
        \textbf{Port} & \textbf{Description} \\
        A1 & Address of the first read register. \\
        A2 & Address of the second read register. \\
        A3 & Address of the write register. \\
        WD & Data to be written to the write register. \\
        RD1 & Data output from the first read register. \\
        RD2 & Data output from the second read register. \\
        WE & Write enable signal. When $1'b1$, data is written to the register A3. \\
        clk & Receives the clock signal (rising edge triggered). \\
        reset & Receives an asynchronous reset signal. \\
        \end{tblr}
        \caption{Register File Module Ports}
        \end{table}
        
        \begin{figure}[H]
            \centering   % Left Image
            \begin{minipage}{0.40\textwidth}
            \centering
            \includegraphics[width=\textwidth]{Register_file_1.png}
            \end{minipage}
            \hfill % Adds horizontal space between the images
            \begin{minipage}{0.40\textwidth}% Right Image
            \centering
            \includegraphics[width=\textwidth]{Register_file_2.png}
            \end{minipage}
            \caption{Register File Internal Structure}
        \end{figure}

        The specific implementation idea of a Register File is relatively straightforward. The \texttt{slt} signal, obtained after decoding the $A3$ input signal via a Decoder, represents the selection signal for each individual register cell. $v$ denotes the value stored in each register cell. The $A1$ and $A2$ input signals are routed through Multiplexers (MUX) to produce $RD1$ and $RD2$, respectively.


        \begin{figure}[H]
            \centering
            \includegraphics[width=0.4\textwidth]{Register_file_3.png}
            \caption{Register Cell Internal Structure}
            \label{fig:RF Module Cell Structure}
        \end{figure}

        For each individual register cell:

        \textbf{clk} receives the clock signal, and \textbf{reset} receives an asynchronous reset signal.

        A \textbf{MUX} selects the content to be written into the register on the rising edge.

        The \textbf{write data (WD)} will be written into the register only if both the \texttt{select signal (slt)} and the \texttt{global write enable (WE)} are at a high level. Otherwise, the existing value will be rewritten.

        
        
        \item \textit{ALU}
        \begin{figure}[H]
            \centering
            \includegraphics[width=0.09\textwidth]{ALU_overall.png}
            \caption{ALU Module Overall Structure}
            \label{fig:ALU Module Overall Structure}
        \end{figure}
        \begin{table}[htbp]
        \centering
        \small
        \begin{tblr}{
            colspec = {l X[j]}, % 'l' for left column auto-width, 'X[j]' for flexible justified column
            hlines, % Horizontal lines
            vlines, % Vertical lines
            row{1} = {font=\bfseries}, % Bold header row
            column{1} = {font=\bfseries}, % Monospace font for port names
        }
        \textbf{Port} & \textbf{Description} \\
        SrcA & First operand input. \\
        SrcB & Second operand input. \\
        ALUcontrol & \parbox{5cm}{Receives ALU operation code. \\ $3'b000$ A \texttt{+} B \\ $3'b001$ A \texttt{-} B \\ $3'b010$ A \texttt{\&} B \\ $3'b011$ A \texttt{|} B \\ $3'b100$ A \texttt{<} B} \\
        ALUResult & ALU output. \\
        Zero & Zero flag indicating if the result is zero. Provide the basis for the judgment of the \texttt{bne} conditional jump\\
        \end{tblr}
        \caption{ALU Module Ports}
        \end{table}

        \begin{figure}[H]
            \centering
            \includegraphics[width=0.7\textwidth]{ALU.png}
            \caption{ALU Module Internal Structure}
            \label{fig:ALU Module Internal Structure}
        \end{figure}

        \item \textit{Extend}
        \begin{figure}[H]
            \centering
            \includegraphics[width=0.3\textwidth]{Extend_overall.png}
            \caption{Extend Module Overall Structure}
            \label{fig:Extend Module Overall Structure}
        \end{figure}

        \begin{table}[htbp]
        \centering
        \begin{tblr}{
            colspec = {l X[j]}, % 'l' for left column auto-width, 'X[j]' for flexible justified column
            hlines, % Horizontal lines
            vlines, % Vertical lines
            row{1} = {font=\bfseries}, % Bold header row
            column{1} = {font=\bfseries}, % Monospace font for port names
        }
        \textbf{Port} & \textbf{Description} \\
        Instr$[31:07]$ & Remove the opcode from the machine code and extract the immediate part \\
        ImmSrc & \parbox{7cm}{Receives the immediate type code. \\ $2'b00$ I1-type \\ $2'b01$ S-type \\ $2'b10$ B-type \\ $2'b11$ J-type} \\
        ImmExt & Output the extended immediate value. \\
        \end{tblr}
        \caption{Extend Module Ports}
        \end{table}

        \begin{figure}[H]
            \centering
            \includegraphics[width=0.8\textwidth]{Extend.png}
            \caption{Extend Module Internal Structure}
            \label{fig:Extend Module Internal Structure}
        \end{figure}

        \begin{table}[htbp]
        \centering
        \begin{tblr}{
            colspec = {l X[j]}, % 'l' for left column auto-width, 'X[j]' for flexible justified column
            hlines, % Horizontal lines
            vlines, % Vertical lines
            row{1} = {font=\bfseries}, % Bold header row
            column{1} = {font=\bfseries}, % Monospace font for port names
        }
        \textbf{Immediate Type} & \textbf{Extraction Method} \\
        I1-type & $20${Instr$[31]$},Instr$[31:20]$ \\
        S-type &  $20${Instr$[31]$},Instr$[31:25]$,Instr$[11:7]$ \\
        B-type &  $20${Instr$[31]$},Instr$[07]$,Instr$[30:25]$,Instr$[11:08]$,$1'b0$ \\
        J-type &  $12${Instr$[31]$},Instr$[19:12]$,Instr$[20]$,Instr$[30:21]$,$1'b0$ \\
        \end{tblr}
        \caption{Immediate Extraction Methods}
        \end{table}
    
        \item \textit{Data Memory}

        The idea is also relatively simple; just use a RAM.

        The address width is 16 bits, and the data width is 32 bits.

        \begin{figure}[H]
            \centering
            \includegraphics[width=0.39\textwidth]{Data_Memory.png}
            \caption{Data Memory Module Implementation}
            \label{fig:DM Module}
        \end{figure}

        \item \textit{Control Unit}
        \begin{figure}[H]
            \centering   % Left Image
            \begin{minipage}{0.3\textwidth}
            \centering
            \includegraphics[width=0.5\textwidth]{Control_Unit_overall.png}
            \caption{Control Unit Overall Structure}
            \end{minipage}
            \hfill % Adds horizontal space between the images
            \begin{minipage}{0.6\textwidth}% Right Image
            \centering
            \includegraphics[width=\textwidth]{Control_Unit_secondary_structure.png}
            \caption{Control Unit Secondary Structure}
            \end{minipage}
            \caption{Control Unit Structure}
        \end{figure}

        \begin{table}[H]
        \centering
        \begin{tblr}{
            colspec = {l X[j]}, % 'l' for left column auto-width, 'X[j]' for flexible justified column
            hlines, % Horizontal lines
            vlines, % Vertical lines
            row{1} = {font=\bfseries}, % Bold header row
            column{1} = {font=\bfseries}, % Monospace font for port names
        }
        \textbf{Port} & \textbf{Description} \\
        opcode & Receives the opcode field from the instruction. \\
        funct3 & Receives the funct3 field from the instruction. \\
        funct7 & Receives the funct7 field from the instruction. \\
        ALUop & Output ALU decoder control signal. \\
        PCSrc & Control signal for PC MUX. \\
        ResultSrc & Control signal for Result MUX. \\
        MenWrite & Data Memory write enable signal. \\
        ALUControl & ALU operation control signal. \\
        ALUSrcB & ALU second operand MUX control signal. \\
        ImmSrc & Immediate extension type control signal. \\
        RegWrite & Register File write enable signal. \\
        \end{tblr}
        \caption{Control Unit Module Ports}
        \end{table}

    
        
        The Control Unit is divided into the \textbf{Main Decoder} and the \textbf{ALU Decoder}.
        
        \begin{figure}[H]
            \centering
            \includegraphics[width=0.75\textwidth]{Main_decoder.png}
            \caption{Main Decoder Internal Structure}
            \label{fig:Control Unit Main Decoder}
        \end{figure}
        The logical approach of the Main Decoder is to first split the opcode$[6:0]$ into a 7-bit binary number, determine the instruction type through \texttt{AND} logic, and then generate the output status signals via \texttt{OR} logic.

        \begin{figure}[H]
            \centering
            \includegraphics[width=0.75\textwidth]{ALU_decoder.png}
            \caption{ALU Decoder Internal Structure}
            \label{fig:Control Unit ALU Decoder}
        \end{figure}

        The logical approach of the ALU Decoder is to first split funct$3$ and funct$7$ into a 3-bit and a 7-bit binary number, respectively. It then determines the exact instruction name by performing \texttt{AND} logic operations on opcode$[5]$, funct7$[5]$, and the ALUOp signal output by the Main Decoder. Finally, it generates the output status signals via \texttt{OR} logic operations.
    \end{enumerate}
\end{itemize}

\subsection{Data Path Explanation}
The data path controls the flow trajectory of data among various modules from the fetch to the execution of each instruction. Controlling the data flow direction according to different instructions is the key to instruction implementation, which is achieved by using Multiplexers (MUX) and their selection signals.

We use 3 MUXes in data path design:
\begin{itemize}  
    \item \textbf{MUX for Next PC selection (PCSrc)}: This MUX selects between the next sequential PC address (PC + 4) and the branch target address (calculated by adding ImmExt to the current PC) as the next PC value. For non-branch instructions, it selects PC + 4; for bne instruction, if the Zero flag is not set, it selects the branch target address.
    \begin{figure}[H]
        \centering
        \includegraphics[width=0.3\textwidth]{PCSrc_MUX.png}
        \caption{MUX for Next PC selection (PCSrc)}
        \label{fig:PCSrc MUX}
    \end{figure}

    PCNext has two input options: pc $+ 4$ and PCTarget. The selection is controlled by the PCSrc signal. For non-branch instructions, PCSrc is $1'b0$, selecting pc $+ 4$ as the next PC value. For the jal instruction, PCSrc is $1'b1$, selecting PCTarget as the next PC value. For the bne instruction, if the Zero flag is not set, PCSrc is $1'b1$, selecting PCTarget as the next PC value.
    \item \textbf{MUX for ALU second operand (ALUSrcB)}: This MUX selects between the second read data from the Register File (RD2) and the extended immediate value (ImmExt) as the second operand for the ALU. For R-type instructions, it selects RD2; for I1-type and S-type instructions, it selects ImmExt.
    \begin{figure}[H]
        \centering
        \includegraphics[width=0.3\textwidth]{ALUSrcB_MUX.png}
        \caption{MUX for ALU second operand (ALUSrcB)}
        \label{fig:ALUSrcB MUX}
    \end{figure}

    ALUSrcB has two input options: RD2 and ImmExt. The selection is controlled by the ALUSrcB signal. For R-type instructions, ALUSrcB is $1'b0$, selecting RD2 as the second operand for the ALU. For I1-type and S-type instructions, ALUSrcB is $1'b1$, selecting ImmExt as the second operand for the ALU. 
    \item \textbf{MUX for Result selection (ResultSrc)}: This MUX selects between the ALU result (ALUResult) and the data read from Data Memory (DMReadData) as the data to be written back to the Register File. For R-type, I1-type, and J-type instructions, it selects ALUResult; for lw instruction, it selects DMReadData.
    \begin{figure}[H]
        \centering
        \includegraphics[width=0.3\textwidth]{ResultSrc_MUX.png}
        \caption{MUX for Result selection (ResultSrc)}
        \label{fig:ResultSrc MUX}
    \end{figure}

    ResultSrc has three input optional: ALUResult, ReadData, and PC $+ 4$. The selection is controlled by the ResultSrc signal. For R-type, I1-type, and S-type instructions, ResultSrc is $2'b00$, selecting ALUResult as the data to be written back to the Register File. For the lw instruction, ResultSrc is $2'b01$, selecting ReadData as the data to be written back to the Register File. For the jal instruction, ResultSrc is $2'b10$, selecting PC $+ 4$ as the data to be written back to the Register File.
\end{itemize}
    In summary, we present the respective control signals corresponding to each instruction, which also serves as an important basis for the CTRL connection.
    \begin{table}[H]
        \centering
        \begin{tblr}{
            colspec = {l X[j]}, % 'l' for left column auto-width, 'X[j]' for flexible justified column
            hlines, % Horizontal lines
            vlines, % Vertical lines
            row{1} = {font=\bfseries}, % Bold header row
            column{1} = {font=\bfseries}, % Monospace font for port names
        }
        \textbf{Instruction} & \textbf{PCSrc} & \textbf{MemWrite} & \textbf{ALUControl} & \textbf{ALUSrcB} & \textbf{ImmSrc} & \textbf{RegWrite} & \textbf{ResultSrc} \\
        I1/Load & $1'b0$ & $1'b0$ & $3'b000$ & $1'b1$ & $2'b00$ & $1'b1$ & $2'b01$ \\ 
        S/Store & $1'b0$ & $1'b1$ & $3'b000$ & $1'b1$ & $2'b01$ & $1'b0$ & -- \\    
        R/OP & $1'b0$ & $1'b0$ & Varies & $1'b0$ & -- & $1'b1$ & $2'b00$ \\
        I1/OP-Imm & $1'b0$ & $1'b0$ & Varies & $1'b1$ & $2'b00$ & $1'b1$ & $2'b00$ \\
        B/Branch & Varies & $1'b0$ & $3'b010$ & $1'b0$ & $2'b10$ & $1'b0$ & -- \\
        J/JAL & $1'b1$ & $1'b0$ & --- & - & $2'b11$ & $1'b1$ & $2'b10$ \\
        \end{tblr}
        \caption{Control Signals for Each Instruction Type}
        \end{table}

\section{Assembler Building}
\begin{lstlisting}[style=mypython, caption=The Instructions Set, label=lst:instrcutions]
INSTRUCTIONS_SET = {
    "add": {"opcode": "0110011", "funct3": "000", "funct7": "0000000", "type": "R"},
    "sub": {"opcode": "0110011", "funct3": "000", "funct7": "0100000", "type": "R"},
    "or": {"opcode": "0110011", "funct3": "110", "funct7": "0000000", "type": "R"},
    "and": {"opcode": "0110011", "funct3": "111", "funct7": "0000000", "type": "R"},
    "slt": {"opcode": "0110011", "funct3": "010", "funct7": "0000000", "type": "R"},
    "slti": {"opcode": "0010011", "funct3": "010", "type": "I1"},
    "lw": {"opcode": "0000011", "funct3": "010", "type": "I1"},
    "addi": {"opcode": "0010011", "funct3": "000", "type": "I1"},
    "ori": {"opcode": "0010011", "funct3": "110", "type": "I1"},
    "andi": {"opcode": "0010011", "funct3": "111", "type": "I1"},
    "sw": {"opcode": "0100011", "funct3": "010", "type": "S"},
    "bne": {"opcode": "1100011", "funct3": "001", "type": "B"},
    "jal": {"opcode": "1101111", "type": "J"},
}
\end{lstlisting}
\begin{lstlisting}[style=mypython, caption=The Assembler Code, label=lst:assembler]
import argparse
from typing import List, Tuple, Optional

import instructions


def preprocess_instruction(line: str) -> Tuple[str, List[str]]:
    """
    Turn an assembly line into (mnemonic, operands[]).

    Supports:
      - commas: addi x1, x0, 1
      - load/store addressing: lw x1, 0(x2)  -> ["x1","0","x2"]
    """
    parts = line.replace(",", "").replace("(", " ").replace(")", "").split()
    if not parts:
        raise ValueError("Empty instruction")
    mnemonic = parts[0]
    operands = parts[1:]
    return mnemonic, operands


def register_to_binary(register: str) -> str:
    if not register.startswith("x"):
        raise ValueError(f"Invalid register name: {register}")
    try:
        reg_num = int(register[1:])
    except Exception:
        raise ValueError(f"Invalid register name: {register}")
    if not (0 <= reg_num <= 31):
        raise ValueError(f"Register out of range x0-x31: {register}")
    return format(reg_num, "05b")


def immediate_to_binary(immediate: str, bits: int) -> str: #please learn: can not read this yet
    imm = int(immediate, 0)  # allow decimal and 0x.. forms
    if imm < 0:
        imm = (1 << bits) + imm
    return format(imm & ((1 << bits) - 1), f"0{bits}b")


def bin_to_hex32(bin_str: str) -> str:
    """32-bit binary string -> 8-hex-digit string (no 0x prefix)."""
    if len(bin_str) != 32:
        raise ValueError(f"Expected 32-bit machine code, got {len(bin_str)} bits")
    return format(int(bin_str, 2), "08x")


def assemble_to_machine_code(assembly_instruction: str) -> dict:
    mnemonic, operand = preprocess_instruction(assembly_instruction)
    inst = instructions.INSTRUCTIONS_SET.get(mnemonic)
    if inst is None:
        raise ValueError(f"Unknown mnemonic: {mnemonic}")

    opcode = inst["opcode"]
    inst_type = inst["type"]

    if inst_type == "R":
        # R-type: mnemonic rd, rs1, rs2
        rd, rs1, rs2 = operand
        machine_code = (
            f"{inst['funct7']}"
            f"{register_to_binary(rs2)}"
            f"{register_to_binary(rs1)}"
            f"{inst['funct3']}"
            f"{register_to_binary(rd)}"
            f"{opcode}"
        )

    elif inst_type == "I1":
        # I-type: addi/andi/ori - rd, rs1, imm
        # lw after preprocess is rd, imm, rs1
        if mnemonic == "lw":
            rd, imm, rs1 = operand
        else:
            rd, rs1, imm = operand
        imm_binary = immediate_to_binary(imm, 12)
        machine_code = (
            f"{imm_binary}"
            f"{register_to_binary(rs1)}"
            f"{inst['funct3']}"
            f"{register_to_binary(rd)}"
            f"{opcode}"
        )


    elif inst_type == "S":
        # S-type: sw rs2, imm(rs1) -> preprocess gives [rs2, imm, rs1]
        rs2, imm, rs1 = operand
        imm_binary = immediate_to_binary(imm, 12)
        machine_code = (
            f"{imm_binary[:7]}"
            f"{register_to_binary(rs2)}"
            f"{register_to_binary(rs1)}"
            f"{inst['funct3']}"
            f"{imm_binary[7:]}"
            f"{opcode}"
        )

    elif inst_type == "B":
        # B-type: bne rs1, rs2, imm
        rs1, rs2, imm = operand
        imm_binary = immediate_to_binary(imm, 13)  # includes sign bit
        # imm[12] imm[10:5] rs2 rs1 funct3 imm[4:1] imm[11] opcode
        machine_code = (
            f"{imm_binary[0]}"
            f"{imm_binary[2:8]}"
            f"{register_to_binary(rs2)}"
            f"{register_to_binary(rs1)}"
            f"{inst['funct3']}"
            f"{imm_binary[8:12]}"
            f"{imm_binary[1]}"
            f"{opcode}"
        )

    elif inst_type == "J":
        # J-type: jal rd, imm
        rd, imm = operand
        imm_binary = immediate_to_binary(imm, 21)
        # imm[20] imm[10:1] imm[11] imm[19:12] rd opcode
        machine_code = (
            f"{imm_binary[0]}"
            f"{imm_binary[10:20]}"
            f"{imm_binary[9]}"
            f"{imm_binary[1:9]}"
            f"{register_to_binary(rd)}"
            f"{opcode}"
        )

    else:
        raise ValueError(f"Unknown instruction type: {inst_type}")

    return {"Assembly": assembly_instruction, "MachineCode": machine_code, "Type": inst_type}


def strip_comment(line: str) -> str:
    # support '#' and ';' comments
    for sep in ("#", ";"):
        if sep in line:
            line = line.split(sep, 1)[0]
    return line.strip()


def assemble_file(
    in_path: str,
    out_path: Optional[str] = None,
    emit_hex: bool = True,
) -> List[str]:
    """
    Assemble a .s file line by line.
    Returns output lines (also writes to out_path if provided).
    """
    output_lines: List[str] = []
    pc = 0  # each instruction 4 bytes (RV32I)
    with open(in_path, "r", encoding="utf-8") as f:
        for lineno, raw in enumerate(f, start=1):
            line = strip_comment(raw)
            if not line:
                continue
            try:
                res = assemble_to_machine_code(line)
                bin_code = res["MachineCode"]
                if emit_hex:
                    hex_code = bin_to_hex32(bin_code)
                    output_lines.append(f" {hex_code}  ")
                else:
                    output_lines.append(f" {bin_code} ")
                pc += 4
            except Exception as e:
                output_lines.append(f"ERROR line {lineno}: {e}    # {raw.rstrip()}")
                # keep going

    if out_path:
        with open(out_path, "w", encoding="utf-8") as wf:
            wf.write("\n".join(output_lines) + "\n")
    return output_lines


def main():
    parser = argparse.ArgumentParser(description="Simple RV32I assembler (subset) - batch mode")
    parser.add_argument("-f", "--file", help="Input assembly file (.s)")
    parser.add_argument("-o", "--out", help="Output text file (optional)")
    parser.add_argument("--bin", action="store_true", help="Emit binary instead of hex")
    args = parser.parse_args()

    if args.file:
        lines = assemble_file(args.file, out_path=args.out, emit_hex=not args.bin)
        print("\n".join(lines))
    else:
        # fallback: single-line interactive mode (original behavior)
        assembly_instruction = input("Enter assembly instruction: ").strip()
        result = assemble_to_machine_code(assembly_instruction)
        print("Assembly Instruction:", result["Assembly"])
        print("Machine Code:", result["MachineCode"])
        print("Instruction Type:", result["Type"])


if __name__ == "__main__":
    main()

\end{lstlisting}
The Assembler converts RISC-V assembly insturctions into 32-bit machine code. It first parses each instruction with \texttt{preprocess\_insturctions} to extract the mnemonic and operands (handling commas, parentheses, and load/store addressing). 
Then \texttt{assemble\_to\_machine\_code} assembles the machine code based on the instruction type: it converts register names to 5-bit binary, converts immediates to the required bit width (using two's complement for negatives),and concatenates fields in the order specified by each instruction format.
Finally, \texttt{assemble\_file} processes assembly files line by line, converts each instruction, and outputs hex or binary while maintaining a PC that increments by 4 bytes per inctruction.



\section{Test Cases and Verification}
\subsection{Test Cases}
We have designed a series of test cases, and judge whether the CPU executes instructions correctly by monitoring the states of the Register File and the Data Memory.
\begin{lstlisting}[style=myriscv, caption=Test Cases, label=lst:testcases]
addi x1, x0, 10  # Initialize x1 with 10
addi x2, x0, -20 # Initialize x2 with -20
addi x3, x0, 5 # Initialize x3 with 5
add x4, x1, x2 # x4 = -10
add x5, x1, x3 # x5 = 15
sub x6, x1, x2 # x6 = 30
sub x7, x1, x3 # x7 = 5
and x8, x1, x2 # x8 = 8
and x9, x1, x3 # x9 = 0
or x10, x1, x2 # x10 = -18
or x11, x1, x3 # x11 = 15
slt x12, x1, x2 # x12 = 0
slt x13, x2, x1 # x13 = 1
slti x14, x1, -8 # x14 = 0
slti x15, x1, 28 # x15 = 1
addi x16, x16, 1 # x16 = 1
addi x17, x0, 3 # x17 = 3
bne x16, x17, -8 # branch to addi x16, x16, 1 if x16 != x17
jal x16, 8 # jump to addi x19, x1, 8
addi x17, x0, 7 # x17 = 7
andi x19, x1, 8 # x19 = 8
andi x20, x2, 18 # x20 = -2
ori x21, x1, 8 # x21 = 10
ori x22, x2, 18 # x22 = -2
sw x19, 0(x5) # Store x19 at address in x5
sw x19, 4(x5) # Store x19 at address in x5 + 4
lw x18, 0(x5) # Load from address in x5 to x18
\end{lstlisting}

\subsection{Verification}
Importing the \texttt{output.txt} file (generated by converting \texttt{testcase.s} via \texttt{Assembler\_batch.py}) into \texttt{Convert.py} will yield a hexadecimal file compatible with Logisim-Evolution v3.0, which can be directly imported into the ROM for automated testing.
\begin{figure}[H]
    \centering
    \begin{minipage}{0.4\textwidth}
        \centering
        \includegraphics[width=\textwidth]{Registerfile_verification.png}
    \end{minipage}
    \hspace{0.1\textwidth}
    \begin{minipage}{0.3\textwidth}
        \centering
        \includegraphics[width=\textwidth]{RAM_verification.png}
    \end{minipage}
    \caption{Verification Sign}
    \label{fig:Verification}
\end{figure}
Testing confirms that our CPU operates normally.
\section{Summary of the project}
This project involved designing and implementing a simple single-cycle RISC-V CPU that supports a subset of the RV32I instruction set. The CPU architecture includes essential components such as the Program Counter (PC), Instruction Memory, Register File, ALU, Data Memory, and Control Unit. Each module was carefully designed to ensure correct functionality and seamless integration within the CPU.

Students have initially learned and understood the RISC-V instruction set and the basic architecture of computers, laying a solid foundation for their subsequent studies.
\section{Team Roles and Responsibilities}
\begin{itemize}
    \item \textbf{Mingxuan Li:} Responsible for assembler development and test case design.
    \item \textbf{Ruiyang Xu:} Responsible for the design and implementation of key modules and the writing of the final report.
    \item \textbf{Kaiwen Yang:} Responsible for datapath implementation and test case design, CPU testing
\end{itemize}

\section{Acknowledgements}
We would like to express our sincere gratitude to our course teaching assistants for their invaluable guidance and support throughout this project. Their expertise and encouragement have been instrumental in helping us navigate the complexities of CPU design and implementation.

\end{document}